\chapter{Conclusion and Future Work}
\label{ch:conclusion}

The dissertation set out to test, on multilingual utility bills, whether direct
vision-mediated extraction by general-purpose LLMs matches or exceeds an
OCR-mediated cascade through the same LLMs, whether the operational trade-off
between the two pipelines splits cleanly, and whether either finding generalises
across LLM providers. The pre-registered hypotheses
(\S\ref{sec:objectives-hypotheses}) committed to specific directional answers
before the runs. The four-cell factorial in
\S\ref{sec:experimental-spec} and the run results in
Chapter~\ref{ch:evaluation} settle each one, and the synthesis of
\S\ref{sec:synthesis} hands six findings forward to this chapter that are
transferable to other document classes and to other LLM stacks.

\section{Transferable findings}
\label{sec:findings}

The following statements summarise what the experiment establishes about the
two pipelines as a category, rather than what was built. Each finding is given
in the form of a transferable claim, paired with a prediction for a new
deployment context the experiment did not cover. The supporting evidence for
each claim sits in the section referenced; this chapter does not repeat the
numbers.

\textbf{Vision-mediated extraction by general-purpose LLMs does not in
itself dominate OCR-mediated extraction on multilingual semi-structured
documents.} The aggregate verdict on H1 (\S\ref{sec:rq1-aggregate}) is that
OCR-mediated extraction outperformed vision-mediated extraction within both
providers tested, by 2.10 percentage points (Sonnet) and 1.02 percentage
points (GPT-4o). The directional prior in the cross-modality literature was
the opposite, and the corpus tested here is the kind of corpus on which the
prior would have predicted the largest vision advantage: short born-digital
multilingual bills with strong layout cues and small extraction schemas.
Applied to a different document class with similar properties (multilingual
invoices, statements of account, expense receipts in a small schema), the
prediction is that an OCR-cascade with a frontier-grade LLM at temperature 0
will at least match a same-LLM vision pipeline, and in some sub-populations
will exceed it. The prediction inverts on long-tail layouts where the field
locations move document-to-document; on those, the per-field flipping
documented in \S\ref{sec:rq1-flipping} predicts that no single-modality
pipeline will be uniformly best.

\textbf{The modality decision is field-level, not pipeline-level.}
Within-document modality flipping affected twelve documents per provider
(\S\ref{sec:rq1-flipping}) and the directions opposed within several of those
documents. A pipeline that selects the modality per field and per condition
would beat the better single-modality pipeline by a margin similar to the
provider-ensemble margin reported in \S\ref{sec:rq3-ensemble}. The
transferable form of this finding is that ``OCR vs vision'' is the wrong
question for any document class with mixed visual and textual cues; the right
question is which fields the document carries primarily as text and which it
carries primarily as layout, and a deployer should expect to route fields
between modalities rather than pick one. Applied to a domain-adjacent
document class such as bank statements, the prediction is that
high-text-density fields (account holder, transaction descriptions) will favour
the OCR-mediated path while layout-heavy fields (current balance, totals
panels) will favour the vision-mediated path.

\textbf{Provider choice is not orthogonal to modality.} The within-provider
modality gap differed between Sonnet and GPT-4o by roughly one percentage point
(\S\ref{sec:rq3-modality-provider}); the providers occupy distinct points in
the hallucination-versus-omission space (\S\ref{sec:rq3-provider}); and the
within-provider modality oracles differ enough that the better single-cell
condition for one provider sits below the better single-cell condition for the
other on several individual fields. Applied to a deployment that does not yet
have a fixed provider, this means the modality choice and the provider choice
must be resolved together, not in sequence: results for one provider on one
modality are a poor predictor of results for the same provider on the other
modality, and an even poorer predictor of results for a different provider on
either.

\textbf{Hallucination and omission are distinct error currencies, and the
choice of provider sets the exchange rate.} Sonnet emitted twenty
hallucinations against twenty omissions; GPT-4o emitted five hallucinations
against thirty-two omissions (\S\ref{sec:rq3-provider}). A downstream system
that punishes wrong values (a payment automation, a fraud check, a tax
filing) inherits a different cost profile from the same accuracy figure
depending on which provider produced it. The transferable claim is that
provider selection on extraction tasks should be conditioned on the
downstream cost asymmetry, not on the headline accuracy alone. The
prediction for medical-form extraction, where a hallucinated field carries
different consequences from a missing field, is that the provider preference
will be the opposite of the one a headline accuracy comparison would suggest.

\textbf{OCR-side failures are not predominantly an OCR problem.} The
four-bucket re-attribution in \S\ref{sec:failure-attribution} showed that
roughly half of the failures the strict diagnosis heuristic attributed to
OCR are actually addressable by improving the LLM, the prompt or the
annotation rather than the preprocessing. The transferable claim is that
``OCR-cascade ceiling'' as it appears in the literature
(\cite{lopresti2009ocr}, \cite{nguyen2021post}) is overstated as a barrier
to LLM-mediated extraction at frontier model scale: post-LLM mis-extraction
is a substantial component of the residual error budget, and remediation
programmes that prioritise OCR-engine swaps over prompt and
instruction-following work will leave half the available accuracy on the
table. Applied to scanned-document corpora, where OCR error is qualitatively
worse than on born-digital input, the same four-bucket reading should be the
first analysis a remediation team produces; the relative size of the
LLM-side bucket may shrink on scans but is unlikely to disappear.

\textbf{The two-page baseline is the right design for any extraction task
whose required fields cluster on early pages.} The two-page input cap
(\S\ref{sec:two-page-cap}) reduced cost by roughly the input-length ratio
without measurable accuracy loss on this corpus, because every required
field already lives within the first two pages of every bill. The
transferable claim is the cost-vs-coverage shape this implies: extraction
cost is dominated by input volume on the vision-mediated path and by
Tesseract latency on the OCR-mediated path, and an early-page bound buys
substantial cost reduction whenever the corpus's information distribution
allows it. Applied to a document class where required fields are sometimes
on later pages (medical referrals, multi-section invoices, structured
contracts), the prediction is that a page-bounded baseline will leave a
small fraction of fields unreachable at zero cost, and that a
header-detection heuristic that selectively expands the input will close
that gap at proportional cost. The cost-vs-coverage curve as a function of
page count is the empirical question this finding most directly motivates.

\section{Limitations}
\label{sec:conclusion-limitations}

Six bounds on external validity are recorded in \S\ref{sec:validity}. Two are
immediate enough to qualify the headline findings here. The corpus is
$n = 40$ active documents, and the per-language and per-utility-type slice
results are descriptive at this size. Differences of a few percentage points
between provider-modality cells should be read as the order-of-magnitude
finding they are, and not as a calibrated effect size. The corpus is also
born-digital only: the OCR error ceiling identified in
\cite{lopresti2009ocr} is not exercised at its hardest, and the
``half of OCR failures are not OCR'' finding is therefore stronger here than
it would be on a scanned-document corpus where the OCR-side bucket would
swell. The other four bounds (Latin-script-only, single-annotator labelling,
single OCR engine, two providers and one model each) qualify specific
findings rather than the headlines, and \S\ref{sec:validity} records each
against the finding it bounds.

The diagnosis-heuristic critique in \S\ref{sec:diagnosis-heuristic} also
bounds itself: the four-bucket re-attribution is a manual reading of the
108 OCR-mode non-correct fields, not a reproducible programme, and a
future replicator who runs the same analysis on a different corpus would
have to redo the manual classification. The strict heuristic remains the
right starting point; the four-bucket reading is the corrective that the
strict heuristic asks of itself when the analyst pauses on each
attribution.

\section{Future work}
\label{sec:future-work}

Six lines of work follow from the findings above. Each is specific, derived
from a section of Chapter~\ref{ch:evaluation}, and stated with the metric
that would settle it.

\textbf{F1: Cost-vs-coverage curve as a function of page count.} The
two-page baseline (\S\ref{sec:two-page-cap}) is right for this corpus
because every required field lives on the first two pages. On a corpus that
exercises later-page targets (statements with detail-page consumption
tables, invoices with totals on a back page, multi-section bills) the
empirical question is how the cost-vs-coverage trade-off shapes against
page count. The experiment to settle it is the present 2 $\times$ 2
factorial replicated at page caps of $\{2, 4, 8, \mathrm{full}\}$ on a
corpus where at least 30\% of required fields live beyond page two; the
metric is per-field accuracy as a function of cap, weighted by the
proportion of fields the cap reaches. The expected result is a piecewise
curve that flattens at the cap that admits all required fields, and a cost
slope in the input-token-length direction that holds across providers.

\textbf{F2: Per-field modality routing.} The within-document modality
flipping in \S\ref{sec:rq1-flipping} and the field-level oracle ceilings
implied by it predict that a router that selects modality per field would
recover most of the ensemble margin reported in
\S\ref{sec:rq3-ensemble}. The router is a small classifier conditioned on
field name, language, and a layout-density signal; the experiment is to
train it on a held-in subset of the present corpus and evaluate it on a
held-out subset. The metric is field-level accuracy against the better
single-modality pipeline for the same provider; the budget is one
classifier per provider. A null result here would falsify the per-field
routing claim and would relocate the ensemble headroom back to provider
ensembling, which is the more expensive option.

\textbf{F3: A second OCR engine.} The OCR-side failure bucket in
\S\ref{sec:failure-attribution} is engine-specific to Tesseract. Replacing
\texttt{TesseractOCREngine} with a drop-in alternative
(Google~Document~AI, AWS~Textract, PaddleOCR) and re-running the four-cell
factorial would isolate which OCR-side failures move when the engine
changes. The metric is the per-field swing on the documents the strict
diagnosis heuristic attributed to OCR; the prediction is that document-AI
class engines will close most of the M3 (locale-aware separator) failures
and a fraction of the M6 (stylised-font misread) failures, while leaving
the LLM-side bucket roughly unchanged.

\textbf{F4: A third LLM provider, and an open-weight provider.} The
provider-effect findings (\S\ref{sec:rq3-provider}) are a two-point
sample. Adding a third frontier provider (Google~Gemini, Mistral~Large) and
an open-weight model (Llama~3.2~Vision, Qwen2-VL) would broaden the
provider axis and would test whether the hallucination-versus-omission
profile extends along a continuum or clusters into discrete provider
personalities. The metric is the per-condition position in the
hallucination-versus-omission plot of \S\ref{sec:rq3-provider}; the
prediction is that frontier models cluster but that the open-weight
profile is meaningfully different on omission rate, because open-weight
models are trained with different RLHF objectives.

\textbf{F5: Scanned-document replication.} The corpus here is born-digital
only. Scanning each bill at $300$~DPI and $150$~DPI and re-running the
four-cell factorial would isolate the scan-quality effect against the
present findings. The metric is the per-cell accuracy delta between
born-digital and scanned input at each DPI; the prediction is that the
OCR-mediated path's accuracy degrades faster than the vision-mediated
path's, that the four-bucket attribution shifts toward the OCR bucket, and
that the modality verdict on H1 inverts on the harder DPI condition.
This is the single experiment that would most directly test the
transferability of the findings here to the document class for which the
literature's vision-prior was built.

\textbf{F6: Prompt sensitivity.} The runs used a single prompt template
(\texttt{prompts/extraction\_v1.txt}) at temperature 0. Three or four
prompt variants per condition, with explicit-locale-format hints, an
explicit-currency-inference clause for ambiguous-symbol jurisdictions, and
a sign-handling clause for credit-bill totals would test whether the
prompt construct issue identified in \S\ref{sec:validity-internal} can be
closed without changing the model. The metric is per-field accuracy on the
fields the prompt clauses target (\texttt{consumption\_amount},
\texttt{currency}, \texttt{total\_amount\_due}); the prediction is that the
sign-handling clause closes most of the credit-bill total losses and that
the explicit-locale clause closes most of the M3 cases at zero additional
inference cost.

\section{Closing remarks}
\label{sec:closing}

The four-cell factorial returns a sharper question than the one it began
with. The literature's framing was modality versus modality, with provider
and prompt held abstract. The findings here suggest that on the document
class tested, the relevant axes are the field-level decision (which
modality wins where), the provider-personality decision (which kind of
failure the downstream system can absorb), and the page-coverage decision
(how much of the document the schema actually requires). None of these is
captured by an aggregate field-level accuracy figure, and the literature
that reports only that figure is therefore answering a question that is
empirically less informative than the questions a deployer needs to ask.
The contribution this dissertation makes is not the verdict on H1; it is
the framing that turns one question into three, and the evidence that the
three are answerable at the cost of the experiment reported here.
